\CGZoneHeader{example-image-b}{ZONA LESTE}
\CGTextSingleColumn{}{Texto descritivo sobre o setor...}

\CGSectorHeader{guide_cyan}{FACE LESTE}
\begin{CGSectorRoutes}
    \CGRoute{Timbu Voador}{VIIb}{25m}{Fixa}{}[! Precisa trocar as proteções ! TIMBU VOADOR 7b | 22m (9+2)]
    \CGRoute{Fuga de Alcatraz}{VIIa}{20m}{Fixa}{}[! Precisa trocar as proteções ! Mesmo começo da Alcatraz. Fica molhada devido as macambiras. FUGA DE ALCATRAZ 7a | 18m (7+2)]
    \CGRoute{Alcatraz}{5° VIIa E1 D1}{60m}{Móvel}{}[! Precisa trocar as proteções ! Interidata pelas macambiras. ALCATRAZ 5º VIIa E1 D1 (55m) A. Jariedson, Marcos Rocha, E. Menger e R. Freitas Eq.: 7 costuras + Jogo de Camalots (segunda enfiada).]
    \CGRoute{Sociedade Alternativa}{VIIa}{15m}{Fixa}{}[SOCIEDADE ALTERNATIVA 7a | 15m (5+2) Dinâmico na saida, fenda de esquerda.]
    \CGRoute{Ecos do Além}{VI}{15m}{Fixa}{}[\CGstar2 ECOS DO ALÉM 6º | 15m (6+2) Inicio com movimento de força esticado para os pequenos. Segue em uma sequencia de movimentos delicados após a terceira chapa, com ótimas agarras para equipar/costurar. FA:]
    \CGRoute{Araramboia}{VI}{15m}{Fixa}{}[\CGstar3 ARARAMBOIA 6º | 15m (6+2) Estilo atletico explosivo.
Começo esticado, segue em Agarras grandes. Final tecnico com regletes medios. FA:]
    \CGRoute{Anaconda}{6° VIIb E2 D1}{60m}{Fixa}{}[\CGstar1 ANACONDA 6º VIIb E2 D1 (50m)]
    \CGRoute{Surucucu}{VIIIa}{20m}{Fixa}{}[SURUCUCU 8a | 20m (7+2) Originalmente, a via era (6+2) e com uso de móvel na saida (compartilhada com a Anaconda). Posteriormente (2026) Júlio Francês adicionou uma chapa na saida (7+2).]
    \CGRoute{Cobra Coral}{VIsup}{30m}{Fixa}{}[COBRA CORAL 6sup | 25m (8+2) Via longa em regletes com lances em barrigas. Exige reistência e boa leitura.]
    \CGRoute{Piriguete}{VIIIa}{25m}{Fixa}{}[PIRIGUETE 8a | 25m (9+2)]
    \CGRoute{Ilhoa}{VIIIa}{25m}{Fixa}{}[ILHOA 8a | 25m (9+2) Compartilha mesma saida da Piriguete]
    \CGRoute{I.A.}{VIsup}{30m}{Fixa}{}[\CGstar3 I.A 6sup | 25m (6+2) Via esportiva mais diferente do setor. 
Começa em escalada em vertical em regletes tecnica até a terceira costura. Segue em uma caverna com lances negativos em agarras enormes. Termina com um mantle ("borda de piscina") em um platô. Parada em ótimo estado com argolas em inox.]
    \CGRoute{Diamante de Mendigo}{4° Vsup E3 D1}{50m}{Móvel}{}[DIAMANTE DE MENDIGO 4º Vsup E3 D1 (50m) Final tomado pela vegetação Inicio em uma fenda (móvel). 
Após segue fixa. Jariedson André, Édervan Menger e Reginaldo Freitas Equipamentos: 
6 Costuras + 
Hexentrics \#.3 e 5 e 
Friends \#.3, 4 e 7]
    \CGRoute{Intrépido}{Vsup}{70m}{Fixa}{}[INTRÉPIDO 5sup (62m) Ary Pacheco Neto Equipamentos: 6 Costuras]
    \CGRoute{Los Manos}{3° Vsup E1 D1}{70m}{Fixa}{}[LOS MANOS 3º Vsup E1 D1 (65m) Alex Patrulha e Jariedson André Equipamentos: 8 Costuras Crux na saida da segunda enfiada.]
    \CGRoute{Aplysia}{3° Vsup E2 D1}{70m}{Fixa}{}[\CGstar1 APLYSIA 3º Vsup E2 D1 (65m) Tem parada intermediaria antes do crux da primeira enfiada. Iraê Terra e Júlio Francês Equipamentos: 11 Costuras]
    \CGRoute{Filhos da Caiada}{V}{30m}{Fixa}{}[\CGstar1 FILHOS DA CAIADA 5º | 25m (5+2)]
    \CGRoute{Dublin}{3° Vsup E3 D1}{50m}{Móvel}{}[DUBLIN 3º Vsup (50m) Ary Pacheco Neto Equipamentos: 14 Costuras, Camalots do \#.3 ao .5 e 4 Cume pelo final da via Aplysia. Parada/Rapel somente em proteção simples]
    \CGRoute{Sapatilha de Hobbit}{VI}{15m}{Fixa}{}[\CGstar1 SAPATILHA DE HOBBIT 6º | 15m (4+2) Parada coberta pela vegetação/macambira. Via em grampos P.]
    \CGRoute{Os Pioneiros}{4° V E2 D1}{60m}{Fixa}{}[OS PIONEIROS 4º V E2 D1 (55m) [Precisa de manutenção] Clodoaldo Dante e João José de Andrade Equipamentos: 8 Costuras e 1 Corda de 60m]
    \CGRoute{Ladrão de Vias}{VI}{30m}{Fixa}{}[\CGstar2 LADRÃO DE VIAS 6º | 30m (12+2)]
    \CGRoute{Wolgrand Imortal}{3° VI E2 D1}{70m}{Fixa}{}[WOLGRAND IMORTAL 3º VI E2 D1 (70m) Ary Pacheco Neto e Rebeca Lins Equipamentos: 8 Costuras Ultimo chapa antes da segunda parada encoberta pela vegetação.]
    \CGRoute{Escopião Martelo}{4° V E2 D1}{60m}{Fixa}{}[ESCORPIÃO MARTELO 4º V E2 D1 (60m) Ary P.Neto, D. Malloy, E. Menger e R. Soares Equipamentos: 12 Costuras Sem parada intermediaria. Descida por trilha ou rapel com 2 cordas de 60m. Ary Pacheco Neto, Denn Malloy, Edervan Menger, Ranieri Soares]
    \CGRoute{25 de Março}{3° IVsup E3 D1}{60m}{Fixa}{}[25 DE MARÇO 3º IVsup E3 D1 (60m) Ary Pacheco Neto Equipamentos: 11 costuras Sem parada intermediaria. Descida por trilha ou rapel com 2 cordas de 60m.]
    \CGRoute{De volta para escola}{III}{8m}{Fixa}{}[DE VOLTA PARA A ESCOLA 3 | 8m (5+2) Escola para Rapel.]
    \CGRoute{Dean Potter}{4° VI E4 D1}{60m}{Fixa}{}[DEAN POTTER 4º VI E4 D1 (60m)]
    \CGRoute{Filipe Carrilho}{3° Vsup E3 D1}{60m}{Fixa}{}[FILIPE CARRILHO 3º V sup (60m) Sem bolts na segunda enfiada. "Necessário 3 removivel bolts"
7 costuras
1 corda de 60m
parada dupla]
    \CGRoute{Ary Filho}{3° IIIsup E2}{60m}{Fixa}{}[\CGstar1 ARY FILHO 3º IIIsup E2 D1 (60m) 8 costuras, 
1 corda de 70m primeira enfiada (7+2) 25m segunda enfiada (7+2) 35m]
    \CGRoute{Genesis}{3° IIIsup E2}{50m}{Fixa}{}[\CGstar1 GENESIS 3º IIIsup E2 D1 (50m)]
    \CGRoute{Malu de Andrade}{3° IVsup E1}{50m}{Fixa}{}[MALU DE ANDRADE 3º IVsup E1 D1 (50m) 7 costuras Via terceiro grau quase vertical com agarras em geral, e com um lance de IVsup em uma barriga bem protegido. Muito vento e difícil comunicação na segunda enfiada.]
    \CGRoute{Itália}{3° VIsup AID:A0}{50m}{Fixa}{}[ITÁLIA 3º VIsup (A0/VIsup) (50m) 13 costuras Lance A0/VIsup é esticado Compartilha parada com a P2 da Malu de Andrade A via conta com 3 proteções em Pitons-Mosquetões grandes não entram. Riccardo Bessio]
    \CGRoute{Miss Simparia}{4° Vsup E1 D1}{50m}{Fixa}{}[MISS SIMPATIA 5º | 50m (11+2) Acesso ao cume Sem parada intermediaria Rapel por outra via Ary Pacheco e Rebeca Lins]
    \CGRoute{Invasoras}{V}{50m}{Fixa}{}[INVASORAS 5º | 50m (11+2) Compartilha proteções com a via Miss Simpatia Não tem parada intermediaria Acesso ao Cume | 
Rapel por outra via]
\end{CGSectorRoutes}

\CGSectorHeader{guide_cyan}{FALÉSIA}
\begin{CGSectorRoutes}
    \CGRoute{Mulambos não tem o que fazer aqui}{VI}{12m}{Fixa}{}[MULAMBOS NÃO TEM O
QUE FAZER AQUI 6º | 12m (4+2) Via mais à esqueda do setor; logo após o cabo de aço.]
    \CGRoute{Adios Amigos}{VIIIa}{}{Fixa}{}[\CGstar3 ADIOS AMIGOS 8a | 20m (7+2) Primeira via e mais clássica do setor Tem uma via inacabada a esqueda dela.]
    \CGRoute{Chapeleiro}{VIIIa}{18m}{Fixa}{}[\CGstar1 CHAPELEIRO 8a | 18m (8+2)]
    \CGRoute{Transmento de Pensaferência}{VIIIb}{18m}{Fixa}{}[\CGstar1 TRANSMENTO DE PENSAFERÊNCIA 8b | 18m (9+2)]
    \CGRoute{A Ressurreição do clip stick}{VIIa}{12m}{Fixa}{}[A RESSURREIÇÃO DO CLIP STICK 7a | 12m (4+2)]
    \CGRoute{Lascosse, Reiosse, Votz}{VIIIc}{18m}{Fixa}{}[LASCOSSE, REIOSSE, VOTZ 8c | 18m (5+2) conquistadores Dante e Menger | nota histórica como "a via mais cara do RN". FA:]
    \CGRoute{Monstros S/A}{IXc}{14m}{Fixa}{}[MONSTROS S/A (9c) | 14m (7) FA:]
    \CGRoute{Retorno de Jedi}{IXa}{14m}{Fixa}{}[\CGstar2 RETORNO DE JEDI 9a | 14m (8+2)]
    \CGRoute{Quebra-Queixo}{VIIb}{14m}{Fixa}{}[QUEBRA-QUEIXO 7a/b | 14m (4) não tem parada dupla]
    \CGRoute{Acapulco Forever}{IXb}{}{Fixa}{}[ACAPULCO FOREVER 9b | 10m (3+2) ! Quebrou um agarrão na saída. Provavelmente subiu o grau. À direita dessa via, tem uma via inacabada.]
    \CGRoute{Triangulo Amoroso}{V}{}{Fixa}{}[\CGstar1 À esquerda da via Gaia Chuva. Conquistada por Tomazini, Rafaela e Gabriel em 2026. No momento da conquista sendo a via mais fácil e acessivel do setor. Saida bem fácil, possibilitando facil proteção até o segundo bolt. Segue com move usando uma agarra destacada (cuidado para não puxar a laca para fora da rocha).]
    \CGRoute{Gaia Chuva}{VIIb}{10m}{Fixa}{}[GAIA CHUVA 7b | 10m (3+2)]
    \CGRoute{Adrenaline}{VIIa}{7m}{Fixa}{}[ADRENALINE 7a | 7m (2+2)]
    \CGRoute{Do Gueto}{VIIb}{7m}{Fixa}{}[DO GUETO 7b | 7m (2+2) Compartilha a parada com a Pais de Gaia]
    \CGRoute{Pais de Gaia}{VIIa}{}{Fixa}{}[PAIS DE GAIA 7a | 7m (2+2)]
    \CGRoute{Delação Premiada}{Vsup}{35m}{Fixa}{}[DELAÇÃO PREMIADA 5sup | 30m (7+2) Recomendado corda de 70m para o rapel]
\end{CGSectorRoutes}

